% TOP_PROBLEM: {"id": 2861, "title": "Kirby Problem 3.63", "category_id": 7, "category": {"id": 7, "name": "topology", "display_name": "Topology", "description": "Properties preserved under continuous deformations.", "slug": "topology", "order_index": 7, "created_at": "2026-07-31T15:26:25.671Z"}, "status": "open"}
% FIRST_POSTED: null
% ENTRY_KIND: statement_only
% Imported from: batch12.tex; UnsolvedMath ID 2861
\documentclass[11pt]{article}
\usepackage[margin=25mm]{geometry}
\usepackage{fontspec}
\setmainfont{Latin Modern Roman}
\usepackage{amsmath,amssymb,mathtools}
\usepackage{unicode-math}
\setmathfont{Latin Modern Math}
\usepackage{xurl,hyperref}
\hypersetup{colorlinks=true,urlcolor=blue}
\setcounter{secnumdepth}{0}
\setlength{\parindent}{0pt}
\setlength{\parskip}{5pt}
\setlength{\emergencystretch}{3em}
\newcommand{\sref}[2]{\hyperlink{source:#1:#2}{[S#2]}}
\newcommand{\cataloguescope}[1]{\paragraph{Recorded target.}#1}
\begin{document}
% UnsolvedMath ID 2861; source catalogue code KP-3.63
\setcounter{section}{2860}
\section{Kirby Problem 3.63}\label{2861}
{\small\sffamily UnsolvedMath ID 2861 · Topology}
\subsection{Definitions and mathematical statement}

\paragraph{Shared notation.}$\mathbb N=\{1,2,\ldots\}$, and $\mathbb Z,\mathbb Q,\mathbb R,\mathbb C$ denote the integers, rationals, reals and complex numbers. Any other conventions are specified locally.


Let $Y$ be a closed oriented hyperbolic $3$-manifold with its curvature-$-1$ metric, and let $s$ be a spin structure, a lift of its oriented orthonormal frame bundle to $\operatorname{Spin}(3)$. The associated complex spinor bundle carries the self-adjoint Dirac operator $D_s$, formed by composing the spin connection with Clifford multiplication. Its real discrete spectrum is counted with multiplicity. For $\operatorname{Re}z$ sufficiently large put
\[
\eta_{D_s}(z)=\sum_{\lambda\in\operatorname{Spec}(D_s)\setminus\{0\}}
\operatorname{sign}(\lambda)|\lambda|^{-z}.
\]
Its meromorphic continuation is regular at $z=0$; the APS Dirac eta invariant is $\eta_{\rm Dirac}(Y,s)=\eta_{D_s}(0)$. This is the unreduced eta invariant, not $(\eta_{D_s}(0)+\dim\ker D_s)/2$.
\paragraph{Statement recorded in the supplied source.}
Give a method to compute $\eta_{\rm Dirac}(Y,s)$ from the hyperbolic manifold and its spin structure, for every closed spin hyperbolic $3$-manifold. The requested quantity is the full real invariant; calculating only its class modulo integers, the signature-operator eta invariant, or examples forced to have zero by an orientation-reversing symmetry does not settle the general computation problem. \sref{2861}{2}
\subsection{Short English statement}
Give a general way to calculate the spectral asymmetry of the Dirac operator for every spin hyperbolic three-manifold.
\subsection{Sources}
\begin{description}
\item[\hypertarget{source:2861:1}{S1}] ULAM.AI and the UnsolvedMath Contributors, \emph{UnsolvedMath: A Curated Collection of Open Mathematics Problems}, record 2861 (KP-3.63). CC BY 4.0. \url{https://huggingface.co/datasets/ulamai/UnsolvedMath}.
\item[\hypertarget{source:2861:2}{S2}] R. İnanç Baykur, Robion C. Kirby and Daniel Ruberman (editors), \emph{K3: A New Problem List in Low-Dimensional Topology}, Mathematical Surveys and Monographs 295 (2026), Problem 3.63, p. 176 and following remarks; original author version. \url{https://math.berkeley.edu/sites/default/files/surv-295-ruberman-watermarked-author-pdf.pdf}.
\end{description}
\label{end:2861}
\end{document}
